\section{Introduction}
Members of the U.S. House of representatives are elected every two years, with a total of $435$ seats apportioned among the states based on population counts from the census~\cite{eckman2019apportionment}.
States that have more than one seat must draw congressional districts and elect one representative per district.
As population distributions shift, district boundary lines are usually redrawn after each census, subject to both federal and state rules.
Since the exact placement of these boundaries can heavily influence election outcomes, the redistricting process is persistently controversial.
Disputes over partisan or racial gerrymandering frequently lead to extensive litigation.
To objectively evaluate whether a redistricting plan has been unfairly manipulated, courts increasingly take into account computational methods~\cite{mattingly2021expert,mattingly2019expert}.

Mathematically, a redistricting plan is a partition of a state's basic geographic or voting units (e.g., precincts) into a fixed number of districts.
These units naturally form a planar graph that encodes geographical contiguity constraints.
A standard way to quantify whether an enacted or proposed redistricting plan is unusual is to compare it against an ensemble of plans that follow similar legal and geographical constraints.
To produce statistically meaningful counterfactuals and hypotheticals, one measures observables of interest~\cite{mattingly2014redistricting,fifield2015new} like voting patterns, demographic distributions, etc., on the ensemble to see whether the enacted plan is typical with respect to the ensemble.
Such ensembles are typically specified via score functions that encode constraints~\cite{mattingly2014redistricting}.
This inspires a large body of work related to drawing ``representative'' redistricting plans ostensibly uniformly at random from all such allowable plans, typically done via MCMC simulations.
Recent algorithmic advances make it feasible to generate large ensembles under increasingly nuanced constraints, for instance, via reversible forest-based merge-split samplers~\cite{autry2020multi,autry2023metropolized} and Cycle Walk~\cite{cyclewalk}.


Despite rapid progress in sampling algorithms, evaluating rare events and exploring complex target measures remains a persistent challenge in the study of redistricting ensembles.
For instance, it is well known that existing samplers exhibit low acceptance rates and slow convergence when the target measure deviates from the tree-count measure~\cite{autry2020multi,cyclewalk}.
To use variance-reduction techniques like stratified sampling, we must first stratify the phase space.
Also known historically in computational chemistry as umbrella sampling~\cite{torrie1977nonphysical}, this approach has become a popular technique for sampling rare events, as well as complex measures~\cite{dinner2020stratification}.
Recent works extend this approach to nonequilibrium sampling via trajectory stratification~\cite{dinner2018trajectory}.
% strata / a set of strata / a stratum
A key prerequisite for stratified sampling is that we need strata that cover the phase space, while maintaining suitable overlaps.
In computational chemistry applications, the underlying phase space is usually an interval, or, more generally a subset of $\mathbb{R}^n$, where there are several canonical ways to construct strata with desired features.
However, it is not immediately clear how to construct such strata over the space of redistricting plans, which is discrete and combinatorial.
Three features make the stratification of redistricting plans particularly challenging: (i) small local boundary moves can produce distinct plans that are, for practical purposes, ``almost the same'', so any useful notion of similarity must capture essential structure while ignoring inconsequential boundary noise; (ii) plans are unordered tuples of districts, so comparing two plans typically requires solving a district-matching problem, which is expensive at scale; and (iii) clustering plans directly can be computationally heavy and may obscure district-level structure, since variation in one district can be largely independent of variation in another.

In this paper, we build a scalable computational pipeline that stratifies redistricting plans. 
We address the aforementioned challenges by first clustering at the district level and then using that structure to build plan-level strata.
This district-level clustering step addresses (i) by grouping together similar districts, and it addresses (iii) by retaining district-level granularity that can be obscured if we directly perform plan-level clustering.
We use a beam-search heuristic to identify, for each plan, a small set of candidate nearest-neighbor words when constructing the strata.
As the first step of the pipeline, we collect districts appearing across the ensemble and cluster them with a hierarchical correlation clustering algorithm \hccultra~\cite{DBLP:journals/corr/abs-2409-01010} under a population-weighted symmetric-difference metric, producing a small ``alphabet'' of representative districts (``letters'').
Each plan is then encoded as an unordered tuple of letters (a ``word'').
Because only a small fraction of all possible words correspond to viable plans, we construct a representative word set near the observed samples from the ensemble via a beam search that expands partial words district-by-district while retaining only the best candidates.
Next, we define a partition of unity (POU) that softly assigns each plan to its nearest words.
In this paper, we provide the framework rather than fixing one specific POU formula, allowing for flexibility when applying this framework to stratified sampling.
In practice, one may need to adjust the weighting according to empirical mixing performance.
Nevertheless, we estimate a stationary distribution over strata and an associated flux matrix with a prescribed POU to demonstrate the plausibility of our method.
We illustrate the pipeline on two real-world examples: Connecticut (small-scale) and North Carolina (large-scale). 
The small-scale Connecticut example provides a better understanding of the words that we construct, since Connecticut only has five districts.
On the other hand, North Carolina has fourteen districts and much more precincts.
We use the North Carolina example to demonstrate the scalability of our pipeline on large-scale datasets.

\medskip
\noindent\textbf{Road map}
%labels compatible with thesis chapter
\begin{itemize}
    \item Section~\ref{sec:chap4 background} provides background on the three ingredients in our pipeline: redistricting ensemble sampling, stratified sampling, and hierarchical clustering.
    \item Section~\ref{sec:chap4 notation_setup} formalizes notation for districts and plans; defines district-level and plan-level distances; introduces letters and words; defines the POU, weights and flux matrix.
    \item Section~\ref{sec:chap4 methodology} details the full pipeline: data generation, hierarchical clustering for building letters, and beam search for candidate words; specifies strategies for choosing alphabet size and word pruning (when overlap is large).
    \item Section~\ref{sec:chap4 numerical_experiments} applies the pipeline to two real-world examples with different sizes (Connecticut and North Carolina).
    \item Section~\ref{sec:chap4 future_directions} discusses next steps, including running full stratified sampling using the strata constructed with the pipeline, and updating the alphabet for new data.
    \item Section~\ref{sec:chap4 conclusion} concludes the paper.
\end{itemize}
